\subsection*{4.5 Proof of Main Theorem}
The proof of our lower bound for the tautology $F^{\mathsf{SEL}}_{n,k}$ with extension axioms $\mathsf{Ext}$ proceeds by choosing a variable in the given refutation $\Pi$ that contributes to a lot of high quadratic degree pairs of terms in $\Pi$. If this variable is Singular, we apply the restriction that sets it to zero. On the other hand, if it is Nonsingular and therefore an extension variable $z$, we first reduce it to depend on a single variable $w^*$ by applying a restriction chosen from Lemma22, and then use a more complicated case analysis (see Lemma30) in order to apply the Split operation from Lemmas23 and25 on $z$. In both of these cases we are able to remove a small fraction of high Quadratic degree terms, and thus after sufficiently many iterations we obtain a refutation of low Quadratic degree. We convert this to a refutation of low (usual) degree using Lemma21, and then substitute for the pigeon variables $\mathbf y$ to select a subset of equations from $F_{n,k}$ that require high degree, obtaining a contradiction.
\subsubsection*{4.5.1 Cleanup operations}
In order to get the contradiction at the end of the above argument, we need to ensure that our process above is always working with a subset of equations of $F_{n,k}$ that are untouched, i.e. unaffected by earlier restrictions to variables. We also need to eliminate any partially assigned pigeons so that we have full choice over the equations we are able to pick at the end. Additionally, a key requirement of the Split lemmas (Lemmas23 and25) is that the variable $z$ we Split on must not appear in any axioms except for one of the form $(z-a)(z-b)=0$, which indicates that it takes two distinct values. In particular, we cannot set $z$ or the underlying variable $w^*$ described above in order to eliminate them from the refutation. This presents us with a unique requirement: for any choice of a variable $w^*\in\mathbf x\cup\mathbf y$, we need to be able to eliminate all axioms containing $w^*$ without actually setting it. We show how to perform these operations by making use of the structure of our tautology $F^{\mathsf{SEL}}_{n,k}$.
We first show how to ``ban'' an equation $E_{b_1\ldots b_{\log m}}$ from $F_{n,k}$ by switching to a set of axioms that prevent any pigeon from being assigned to $b_1\ldots b_{\log m}$.
\paragraph{Lemma26.} Let $\Pi$ be a refutation of $F^{\mathsf{SEL}}_{n,k}|_{\rho}$ for some restriction $\rho$ and let $(\mathbf y_i\neq b_1\ldots b_{\log m})\cdot E_{b_1\ldots b_{\log m}}=0$ be one of its axioms. Then there exists another valid refutation $\Pi'$ with the latter axiom replaced by the axiom $(\mathbf y_i\neq b_1\ldots b_{\log m})\equiv\prod_{b_j=1}y_{i,j}\prod_{b_j=0}\bar y_{i,j}$, such that the quadratic degree of $\Pi'$ is at most that of $\Pi$.
\begin{proof}
Note that the axiom $(\mathbf y_i\neq b_1\ldots b_{\log m})\cdot E_{b_1\ldots b_{\log m}}=0$ can be derived from the axiom $(\mathbf y_i\neq b_1\ldots b_{\log m})\equiv\prod_{b_j=1}y_{i,j}\prod_{b_j=0}\bar y_{i,j}$ by multiplying by the polynomial $E_{b_1\ldots b_{\log m}}$. Since this derivation involves only singular variables, the degree can never drop and therefore the quadratic degree of this derivation is at most that of the final polynomial. We construct $\Pi'$ as follows. We first derive the former axiom from the latter in $\Pi'$. Besides this derivation, $\Pi'$ involves the same steps as $\Pi$.
\end{proof}
\paragraph{Definition27.} An equation $E_{b_1\ldots b_{\log m}}$ is said to be banned when the previous lemma is applied repeatedly to eliminate all occurrences of it from the axioms.
\paragraph{Definition28.} A clean version of $F^{\mathsf{SEL}}_{n,k}$ is any subset of axioms of $F^{\mathsf{SEL}}_{n,k}$ along with axioms that ban some subset of equations of the form $E_{b_1\ldots b_{\log m}}$.
\paragraph{4.5.1.1 Cleanup($\rho$).}
We now describe how to perform the cleanup operations, which we collectively call $\mathsf{Cleanup}(\rho)$, that takes as input an ``unclean'' version of $F^{\mathsf{SEL}}_{n,k}$ derived by applying a restriction $\rho$ to a clean version, and outputs another clean version that is in some sense a subset of the input. Suppose that we are given a restriction $\rho$ that has been applied to a clean version of $F^{\mathsf{SEL}}_{n,k}$, with a variable $w^*\in\rho$ possibly set to $\star$, indicating that it must remain unset. To eliminate an axiom that has been affected by a $\mathbf x$ variable in $\rho$ not set to $\star$, we simply obtain the refutation that bans the corresponding equation $E_{b_1\ldots b_{\log m}}$ as described in the above lemma. Note that since we are eliminating the axiom without setting any variables in it, we can also do this in case our variable $w^*\in\mathbf x$. Suppose that $y_{ij}$ is a $\mathbf y$ variable in $\rho$ not set to $\star$. We first note that any axiom that contained $y_{ij}$ before the application of $\rho$ contains all the variables $y_{i1}\ldots y_{i\log m}$ corresponding to the $i$th pigeon $\mathbf y_i$. We first make sure that this $i$th pigeon does not contain our variable $w^*$ that must remain unset. If it doesn't, we proceed as follows. We set all the other variables in this pigeon to select some equation $E_{b_1\ldots b_{\log m}}$ that has not been banned. Such an equation exists provided that the number of banned equations so far is bounded, and the size of the restriction $\rho$ is also bounded (we formalize this in the lemma below). We then apply an additional restriction to the $\mathbf x$ variables that satisfies this equation $E_{b_1\ldots b_{\log m}}$ picked above. We then ban all the equations affected by this additional restriction, like we did above for the part of $\rho$ containing $\mathbf x$ variables. This eliminates the pigeon $\mathbf y_i$. We are left with the case where our variable $w^*$ belongs to some pigeon $\mathbf y_j$. We set all the variables in the pigeon $\mathbf y_j$ except for $w^*$, such that neither of the two equations $E_{b_1\ldots b_{\log m}}$ and $E_{b'_1\ldots b'_{\log m}}$ that would be selected if $w^*$ is set to zero or one are banned (again, these exist under the same conditions as above). We then proceed as before, i.e. apply an additional restriction to satisfy both these equations, and then ban any other equations that have been affected by this additional restriction. With this we have eliminated the axioms of pigeon $\mathbf y_j$ which select an equation, but we are still left with the axioms that prevent $\mathbf y_j$ from colliding with any other pigeon, which are now of the form $w^*\cdot(\mathbf y_{j'}\neq b_1\ldots b_{\log m})$ and $w^*\cdot(\mathbf y_{j'}\neq b'_1\ldots b'_{\log m})$ indicating that any pigeon $\mathbf y_{j'}$ distinct from $\mathbf y_j$ must not be mapped to the equations $E_{b_1\ldots b_{\log m}}$ and $E_{b'_1\ldots b'_{\log m}}$ if one of them is selected by setting $w^*$ to zero or one. To remove the latter axioms we do something similar to the process of banning an equation, where we simply replace these axioms by the axioms $(\mathbf y_{j'}\neq b_1\ldots b_{\log m})$ and $(\mathbf y_{j'}\neq b'_1\ldots b'_{\log m})$, effectively banning the equations $E_{b_1\ldots b_{\log m}}$ and $E_{b'_1\ldots b'_{\log m}}$ for the remaining pigeons.
\paragraph{4.5.1.2 Correctness of Cleanup($\rho$).}
We note that the above cleanup operations over $\mathbf y$ variables terminate successfully only when there are enough equations that have not been banned by prior calls to cleanup, and also the size of the restriction $\rho$ is bounded. We make this formal by the below lemma.
\paragraph{Lemma29 (Correctness of Cleanup($\rho$)).} Let $\rho$ be a restriction of size $\kappa$. If the number of banned equations (from previous calls to Cleanup) is $\ll m/2^{\kappa}$, then $\mathsf{Cleanup}(\rho)$ terminates correctly. Moreover, it bans at most $O(\kappa)$ additional equations and removes at most $O(\kappa)$ pigeons in its run.
\begin{proof}
In $\mathsf{Cleanup}(\rho)$, note that we can remove the axioms that contain $\mathbf x$ variables unconditionally. When we remove a pigeon $\mathbf y_i=y_{i1}\ldots y_{i\log m}$, we rely on having an equation it can be set to that is not already banned. Since the size of $\rho$ is bounded by $\kappa$, note that at most $\kappa$ variables from $y_{i1}\ldots y_{i\log m}$ can be set by $\rho$. Therefore there are at least $\log m-\kappa$ of them unset, corresponding to selecting $m/2^{\kappa}$ many equations. Since we assume that the number of banned equations is much less than this, we can always find one that is not banned to assign this pigeon to.
We now count the number of new equations banned and the number of pigeons removed by this call to $\mathsf{Cleanup}(\rho)$. Since each $\mathbf x$ variable appears in a constant number of equations, the number of equations we ban while processing it is a constant. When we process a $\mathbf y$ variable, we pick and satisfy an equation, and ban all other equations affected in the process. Since every equation also contains a constant number of variables, satisfying it affects only a constant number of other equations. Therefore, for every variable we process we ban only a constant number of equations, and thus the total number of equations banned is $O(\kappa)$. We remove only those pigeons with a variable in $\rho$, so this is also bounded by $O(\kappa)$.
\end{proof}
\subsubsection*{4.5.2 The Main Theorem}
We need first the following key lemma that shows how to apply the Split operation to reduce high quadratic degree terms.
\paragraph{Algorithm1: Algorithm for Lemma30.}
\textbf{Input:} A refutation $\Pi$, and a nonsingular variable $z$ with extension axiom $z-Q=0$ satisfying the pre-conditions of Lemma30.
\textbf{Output:} A refutation $\Pi'$ satisfying post-conditions of Lemma30.
\begin{enumerate}[label=\arabic*,leftmargin=*]
\item Let $\ell_0<\ell_1$ be such that $|H_d(\Pi,z,\ell_0,\ell_1)|\geq|H_d(\Pi,z)|/p^2$.
\item Apply Lemma22 with $l=\ell_1-\ell_0$ to obtain $\delta,w^*,a,b$ satisfying post-conditions of Lemma22.
\item $\Pi=\Pi|_{\delta}$ (and in particular $z-Q|_{\delta}=z-(b-a)w^*-a$).
\item $\mathsf{Cleanup}(\delta\cup\{w^*=\star\})$ (Cleanup axioms affected by $\delta$ and remove $w^*$ from all axioms other than $z-(b-a)w^*-a$ while keeping it alive.)
\item \textbf{if} $w^*$ contributes to $\geq\epsilon/4p^2$ fraction of pairs in $H_d(\Pi)$ \textbf{then}
\item $\Pi=\Pi|_{w^*=0}$
\item \textbf{end}
\item \textbf{else}
\item Apply the substitution $(z-a)/(b-a)$ for $w^*$ in $\Pi$.
\item Let $\Pi'=\Split_{z,l_0,l_1}(\Pi)$.
\item \textbf{end}
\end{enumerate}
